%! AP Statistics — Unit 3, Lesson 3.7 — Homework KEY
\documentclass[10pt]{article}
\usepackage{apstats-article}
\usepackage{apstats-key}
\usepackage{grading}
\gradingassignment{apstats-u03-l07-hw}
\gradingcourse{advanced_statistics}

\begin{document}

\pageheader{Unit 3, Lesson 3.7}{Homework — KEY}

\vspace{0.1in}

\begin{practicebox}
\small
\textbf{1.}

Rand. assignment: \ans{No} \quad Rand. selection: \ans{Yes}\\[3pt]
Causation: \ans{No} \quad Generalize to: \ans{All registered voters nationally}\gradedkey{\mbox{Assign.~\gradedblank[id=hw1a, type=short, points=1, answer=no, check=text]{0.33in}{0.17in}\enspace Sel.~\gradedblank[id=hw1b, type=short, points=1, answer=yes, check=text]{0.6in}{0.17in}\enspace Cause~\gradedblank[id=hw1c, type=short, points=1, answer=no, check=text]{0.33in}{0.17in}}}\gradedkey{\gradedblank[id=hw1d, type=short, points=1, answer={all national registered voters}, check=text, display={\bfseries all national registered voters}]{3.6in}{0.3in}}\vspace{5pt}
\ans{Because no random assignment was used, we cannot conclude that social media use causes more polarized views — confounding variables (e.g., age, news consumption) may explain the association. Because a random sample was used, the finding can be generalized to all national registered voters.}\gradedkeyline{\gradedblank[id=hw1e, type=rubric, points=2, display={No random assignment: can't say social media causes polarization (confounding). Random sample: generalize to all national registered voters.}]{5.3in}{\gradedlinesheight{2}}}

\vspace{6pt}
\textbf{2.}

Rand. assignment: \ans{Yes} \quad Rand. selection: \ans{No (volunteers)}\\[3pt]
Causation: \ans{Yes} \quad Generalize to: \ans{Similar volunteers; not all students}\gradedkey{\mbox{Assign.~\gradedblank[id=hw2a, type=short, points=1, answer=yes, check=text]{0.33in}{0.17in}\enspace Sel.~\gradedblank[id=hw2b, type=short, points=1, answer=no, check=text]{0.6in}{0.17in}\enspace Cause~\gradedblank[id=hw2c, type=short, points=1, answer=yes, check=text]{0.33in}{0.17in}}}\gradedkey{\gradedblank[id=hw2d, type=short, points=1, answer={similar volunteers only}, check=text, display={\bfseries similar volunteers only}]{3.6in}{0.3in}}\vspace{5pt}
\ans{Because random assignment was used, the active-learning classroom caused the significantly higher exam scores. Because the participants were volunteers (not randomly selected), we cannot generalize beyond students similar to these volunteers.}\gradedkeyline{\gradedblank[id=hw2e, type=rubric, points=2, display={Random assignment: active learning caused the higher scores. Volunteers: generalize only to students like them.}]{5.3in}{\gradedlinesheight{2}}}

\tcbbreak
\vspace{6pt}
\textbf{3.}

Rand. assignment: \ans{Yes} \quad Rand. selection: \ans{Yes}\\[3pt]
Causation: \ans{Yes} \quad Generalize to: \ans{Patients at these 10 hospitals (or similar hospital patients)}\gradedkey{\mbox{Assign.~\gradedblank[id=hw3a, type=short, points=1, answer=yes, check=text]{0.33in}{0.17in}\enspace Sel.~\gradedblank[id=hw3b, type=short, points=1, answer=yes, check=text]{0.6in}{0.17in}\enspace Cause~\gradedblank[id=hw3c, type=short, points=1, answer=yes, check=text]{0.33in}{0.17in}}}\gradedkey{\gradedblank[id=hw3d, type=short, points=1, answer={patients at these 10 hospitals}, check=text, display={\bfseries patients at these 10 hospitals}]{3.6in}{0.3in}}\vspace{5pt}
\ans{Because random assignment was used, the new drug caused the statistically significant improvement. Because patients were randomly selected from these hospitals, we can generalize the finding to the patient population at these hospitals.}\gradedkeyline{\gradedblank[id=hw3e, type=rubric, points=2, display={Random assignment: the drug caused the improvement. Random selection: generalize to patients at these 10 hospitals.}]{5.3in}{\gradedlinesheight{3}}}

\vspace{6pt}
\textbf{4.} \ans{Statistical significance means the observed result is too large to be explained by chance alone. Practical importance (effect size) refers to whether the difference is large enough to matter in real life. Example: A study finds that a new drug reduces blood pressure by 1 mm Hg (statistically significant with n = 100,000). This is statistically significant but probably not clinically important — a 1 mm Hg difference has little real-world impact on health outcomes.}\gradedkeyline{\gradedblank[id=hw4, type=rubric, points=2, display={Significance: too large to be chance. Practical importance: large enough to matter. E.g., a 1 mm Hg drop with n = 100,000 is significant but not clinically important.}]{6.1in}{\gradedlinesheight{3}}}
\end{practicebox}

\end{document}
